\section{Conclusion}
\label{sec:conclusion}

Case volume has different statistical value for different inferential targets. Under fixed, known reporting and exogenous exposure, counting error decreases for both the current reproduction number and the realised mean of observed periods. Estimating a stationary process mean additionally requires sampling periods whose transmission levels vary; dependence determines the remaining uncertainty. Unobserved reporting volatility contributes error even for realised targets, and count moments generally cannot separate it from transmission variation. Apparent weekly admission statistics measure deviations around a chosen centre rather than error against known transmission truth. Curve comparison, thinning, state-intercept fits and conditional calibration each address part of this observation problem, without identifying a universal case threshold. Surveillance design should specify the target, observation mechanism and temporal scale before using case volume to quantify precision or allocate resources.
